%%%%%%%%%%%%%%%%%%%%%%%%%%%%%%%%%%%%%%%%%%%%%%%%%%%%%%%%%%%%%
% ==================== 模型的评价、改进与推广 ====================
% BZD数模社制作，本文档为论文模板；如需详细说明，可见论文模板完整版。
% ✓ 自查清单：
%   □ 是否概括四问模型、输入输出关系和主要结论？
%   □ 每项优点是否对应实际做法或量化证据？
%   □ 每项不足是否写清来源、影响和适用边界？
%   □ 改进措施是否逐项对应前述不足？
%   □ 推广是否说明需要重新标定的变量、参数和约束？
%   □ 是否避免重复前文的模型检验与灵敏度分析？

\section{模型的评价、改进与推广}

本文围绕【研究问题】，针对四个子问题分别采用【各问主要模型或方法】完成
【问题一目标】、【问题二目标】、【问题三目标】和【问题四目标】。
其中，问题一得到的【结果】为问题二提供【输入或参数】，问题二与问题三的【中间结果】
进一步进入问题四的【目标函数、约束或综合决策】。结合前文真实求解结果和检验结论，
现对模型的优点、不足、改进方向及推广价值进行总结。

\subsection{模型的优点}

\begin{enumerate}
  \item 模型具有较强针对性。针对【题目特点】，在【基础模型或理论】中引入【变量、机制或约束】，使模型能够描述【实际关系】，并由【真实结果或指标】提供证据。
  \item 模型结构较为完整。通过【目标函数、状态方程、约束体系或求解策略】统一刻画【核心矛盾】，其中【量化指标或验证结果】达到【真实数值或判定标准】。
  \item 各问题衔接清晰。前序问题输出的【参数、特征或方案】真实进入后续模型，全文符号、单位和评价口径保持一致。
  \item 结果具有【可解释性、稳定性或计算效率】。由【对比实验、误差检验或稳健性分析】可知，【填写有证据支持的具体优点】。
\end{enumerate}

\subsection{模型的不足}

\begin{enumerate}
  \item 受【数据来源、样本范围或统计口径】限制，模型尚未充分考虑【缺失因素】，可能使【具体结果】在【适用边界】之外产生偏差。
  \item 为保证模型可解，对【现实过程】作出【具体假设或简化】，因此模型在【极端工况或特殊情景】下的解释能力受到限制。
  \item 【参数估计、算法规模或随机求解过程】仍存在【具体局限】，可能影响【精度、求解时间或决策稳定性】。
\end{enumerate}

\subsection{模型的改进}

针对上述不足，可从以下方面改进：

\begin{enumerate}
  \item 针对数据局限，补充【数据类型】并扩大【时间、空间或样本】范围，统一【统计口径】，重新估计【相关参数】。
  \item 针对模型简化，引入【新增变量、动态机制或现实约束】，并通过【检验方法】评价改进前后的差异。
  \item 针对计算与不确定性问题，采用【改进算法、鲁棒优化或不确定性建模方法】，比较【精度、稳定性与计算成本】后确定改进方案。
\end{enumerate}

\subsection{模型的推广}

本文模型中的【可迁移结构、变量关系或求解思想】可推广至【相近对象或应用场景】。
推广时应重新收集【新场景数据】，标定【关键参数】，替换【场景相关变量】，并增加
【新的资源、边界或业务约束】；随后采用【验证方法】确认模型在新场景中的适用性，
不得直接照搬本文参数、结果或结论。
